\documentclass[11pt,a4paper]{article}

\usepackage{indentfirst}
\usepackage{amssymb}
\usepackage{subcaption}
\usepackage{graphicx}
\usepackage{booktabs}
\usepackage{longtable}
\usepackage{fancyhdr}
\usepackage{xeCJK}
\usepackage{amsmath}
\usepackage{amssymb}
\usepackage{ulem}
\usepackage{xcolor}
\usepackage{fancyvrb}
\usepackage{listings}
\usepackage{soul}
\usepackage{hyperref}
	\lstdefinestyle{C}{
   		language=C, 
   		basicstyle=\ttfamily\bfseries,
    	numbers=left, 
    	numbersep=5pt,
    	tabsize=4,
    	frame=single,
   	 	commentstyle=\itshape\color{brown},
    	keywordstyle=\bfseries\color{blue},
   	 	deletekeywords={define},
    	morekeywords={NULL,bool}
	}	


\setCJKmainfont{Noto Sans Mono CJK TC}
 
\voffset -20pt
\textwidth 410pt
\textheight 650pt
\oddsidemargin 20pt
\newcommand{\XOR}{\otimes}
\linespread{1.2}\selectfont
\pagestyle{fancy}
\lhead{宜蘭高中學科能力競賽測機試題}

\begin{document}

\begin{center}
\section*{A. 簡單題}
\end{center}

\section*{Description}

這是一個相當簡單的題目，計算一些加法就好

\section*{Input}
第一行為一個正整數 $t$，代表子測資數量。

接下來$t$行，每行有兩個正整數$a, b$，保證$|a|, |b| \leq 10^{100}$


\section*{Output}

對於每一組測資，輸出一行，為 $a+b$ 的結果。

\section*{Sample}
\begin{longtable}[!h]{|p{0.5\textwidth}|p{0.5\textwidth}|}
\hline
\textbf {Input}	& \textbf {Output} \\
\hline
\parbox[t]{0.5\textwidth} % sample 1
{ \tt
5\\
1 2\\
3 4\\
5 6\\
7 8\\
9 10\\
} &
\parbox[t]{0.5\textwidth}
{ \tt
3\\
7\\
11\\
15\\
19\\
} \\
\hline
\end{longtable}


\section*{Subtasks}

本題共有 $3$ 個子任務，滿分為 $100$ 分，通過某子任務即可拿到該子任務分數。

本題子任務:

\begin{center}
    \begin{tabular}{||c c c||}
    \hline
    子任務 & 條件限制 & 分數 \\
    \hline
    $1$ & 範例測資 & $0$ \\
    \hline
    $2$ & $-100 \leq a, b \leq 100$ & $30$  \\
    \hline
    $3$ & 題目範圍限制 & $70$ \\
    \hline
    \end{tabular}
\end{center}

\end{document}

